\documentclass[11pt,a4paper]{article}

% ---------------------------------------------------------------- langue & fontes
\usepackage{fontspec}
\usepackage[french]{babel}
\frenchsetup{SmallCapsFigTabCaptions=false,ThinSpaceInFrenchNumbers=true}
\setmainfont{LibertinusSerif-Regular.otf}[
  BoldFont=LibertinusSerif-Bold.otf, ItalicFont=LibertinusSerif-Italic.otf,
  BoldItalicFont=LibertinusSerif-BoldItalic.otf, Numbers=OldStyle]
\setsansfont{Inter-Regular.otf}[
  BoldFont=Inter-Bold.otf, ItalicFont=Inter-Italic.otf,
  BoldItalicFont=Inter-BoldItalic.otf, Scale=0.94]
\setmonofont{FiraMono-Regular.otf}[BoldFont=FiraMono-Bold.otf, Scale=0.84]
\newfontfamily\tabfont{Inter-Regular.otf}[BoldFont=Inter-Bold.otf,
  ItalicFont=Inter-Italic.otf, BoldItalicFont=Inter-BoldItalic.otf,
  Scale=0.94, Numbers=Lining]

% ---------------------------------------------------------------- mise en page
\usepackage[a4paper,top=2.5cm,bottom=2.3cm,left=2.1cm,right=2.1cm,
            headheight=17pt,headsep=13pt,footskip=20pt]{geometry}
\usepackage{microtype}
\usepackage{xcolor}
\usepackage{booktabs,tabularx,longtable,array,colortbl,multirow}
\usepackage{titlesec,titletoc}
\usepackage{fancyhdr}
\usepackage{enumitem}
\usepackage[most]{tcolorbox}
\usepackage{fancyvrb}
\usepackage{pifont}
\usepackage{fontawesome5}
\usepackage{needspace}
\usepackage{ragged2e}
\usepackage[hidelinks]{hyperref}

% ---------------------------------------------------------------- palette
\definecolor{encre}   {HTML}{16181D}
\definecolor{accent}  {HTML}{1F4E5F}
\definecolor{accent2} {HTML}{8C4A2F}
\definecolor{gris}    {HTML}{6B7280}
\definecolor{filet}   {HTML}{D5DADE}
\definecolor{fondclair}{HTML}{F3F6F7}
\definecolor{fondcode}{HTML}{F7F7F4}
\definecolor{vert}    {HTML}{17693F}
\definecolor{rouge}   {HTML}{9B2C2C}
\definecolor{ambre}   {HTML}{A86A12}
\color{encre}

\hypersetup{colorlinks=true,linkcolor=accent,urlcolor=accent,
  pdftitle={Reunion de cadrage - Projet NLP 2},
  pdfsubject={Jalon J2 - choix des cas d'usage, cadrage, outils novateurs},
  pdfauthor={Equipe projet NLP 2 - M2 Data et IA, FGES}}

% ---------------------------------------------------------------- symboles
\newcommand{\okmark}   {\textcolor{vert}{\ding{51}}}
\newcommand{\komark}   {\textcolor{rouge}{\ding{55}}}
\newcommand{\warnmark} {\textcolor{ambre}{\faExclamationTriangle}}
\newcommand{\checkbox} {\textcolor{gris}{\ding{113}}}
\newcommand{\cutmark}  {\textcolor{accent2}{\ding{34}}}
\newcommand{\starmark} {\textcolor{ambre}{\ding{72}}}
\newcommand{\wipmark}  {\textcolor{ambre}{\faTools}}
\newcommand{\code}[1]  {{\ttfamily #1}}
\newcommand{\rulebreak}{}

% ---------------------------------------------------------------- titres
\titleformat{\section}[hang]
  {\sffamily\bfseries\LARGE\color{accent}}{\thesection}{13pt}{}
  [\vspace{2pt}{\color{filet}\titlerule[0.9pt]}]
\titlespacing*{\section}{0pt}{24pt plus 4pt minus 2pt}{11pt}
\titleformat{\subsection}
  {\sffamily\bfseries\large\color{encre}}{\thesubsection}{9pt}{}
\titlespacing*{\subsection}{0pt}{16pt plus 3pt minus 1pt}{6pt}
\titleformat{\subsubsection}
  {\sffamily\bfseries\normalsize\color{accent2}}{}{0pt}{}
\titlespacing*{\subsubsection}{0pt}{12pt plus 2pt minus 1pt}{4pt}
\setcounter{secnumdepth}{2}
\setcounter{tocdepth}{2}

% ---------------------------------------------------------------- en-têtes
\pagestyle{fancy}\fancyhf{}
% en-tête courant, redéfini par build.py selon le document
\newcommand{\entetecourant}{Projet NLP 2}
\fancyhead[L]{\sffamily\scriptsize\color{gris}\entetecourant}
\fancyhead[R]{\sffamily\scriptsize\color{gris}\thepage}
\renewcommand{\headrulewidth}{0pt}
\renewcommand{\headrule}{\vspace{-4pt}{\color{filet}\rule{\headwidth}{0.4pt}}}
\fancypagestyle{plain}{\fancyhf{}\renewcommand{\headrule}{}%
  \fancyfoot[C]{\sffamily\scriptsize\color{gris}\thepage}}

% ---------------------------------------------------------------- paragraphes & listes
\setlength{\parindent}{0pt}
\setlength{\parskip}{5.5pt plus 1.5pt minus 0.5pt}
\linespread{1.05}
\raggedbottom
\widowpenalty=10000 \clubpenalty=10000
\setlist[itemize]{leftmargin=1.15em,itemsep=2pt,topsep=3pt,parsep=0pt,label=\textcolor{accent}{\textbullet}}
\setlist[enumerate]{leftmargin=1.5em,itemsep=2.5pt,topsep=3pt,parsep=0pt,
  label=\textcolor{accent}{\sffamily\bfseries\arabic*.}}

% ---------------------------------------------------------------- tableaux
% \hspace{0pt} : glue initiale, sans quoi TeX refuse de couper le premier
% mot de la cellule — cause de débordements en colonne étroite.
\newcolumntype{L}[1]{>{\raggedright\arraybackslash
  \rightskip=0pt plus 1fil\relax\hspace{0pt}}p{#1}}
% \hspace{0pt} : sans cette glue, le whatsit de \color empêche TeX
% de couper le premier mot de la cellule (débordements en en-tête).
\newcommand{\thead}[1]{{\tabfont\bfseries\color{accent}\hspace{0pt}#1}}
\newenvironment{tabletype}
  {\par\addvspace{7pt}\begingroup
   \tabfont\setlength{\tabcolsep}{4pt}\renewcommand{\arraystretch}{1.26}
   \setlength{\LTleft}{0pt}\setlength{\LTright}{0pt}
   \setlength{\extrarowheight}{1pt}
   % élasticité infinie à droite : TeX ne coupe plus un mot pour « remplir »,
   % mais coupe encore lorsque c'est le seul moyen d'éviter un débordement.
   % césure interdite : les largeurs minimales sont mesurées, chaque mot tient.
   % Un débordement, s'il en restait un, serait signalé par le journal.
   \hyphenpenalty=9000\exhyphenpenalty=9000
   \emergencystretch=0pt\hbadness=10000\relax}
  {\endgroup\par\addvspace{9pt}}
\setlength{\aboverulesep}{0.5pt}\setlength{\belowrulesep}{0.6pt}
\arrayrulecolor{filet}
\setlength{\heavyrulewidth}{0.9pt}\setlength{\lightrulewidth}{0.45pt}

% ---------------------------------------------------------------- encadrés
\newtcolorbox{callout}{enhanced,breakable,
  colback=fondclair,colframe=accent,boxrule=0pt,leftrule=2.6pt,
  sharp corners,left=11pt,right=11pt,top=8pt,bottom=8pt,
  before skip=10pt,after skip=11pt,parbox=false}
\newtcolorbox{codebox}{enhanced,breakable,
  colback=fondcode,colframe=filet,boxrule=0.5pt,arc=2pt,
  left=9pt,right=9pt,top=7pt,bottom=7pt,
  before skip=10pt,after skip=11pt}

% ---------------------------------------------------------------- sommaire
\renewcommand{\contentsname}{Sommaire}
\titlecontents{section}[0pt]
  {\addvspace{4pt}\sffamily\bfseries\color{encre}}
  {\contentslabel[\thecontentslabel]{1.6em}}{}
  {\titlerule*[0.7pc]{\color{filet}.}\contentspage}
\titlecontents{subsection}[1.9em]
  {\sffamily\small\color{gris}}
  {\contentslabel[\thecontentslabel]{2.3em}}{}
  {\titlerule*[0.7pc]{\color{filet}.}\contentspage}

% ---------------------------------------------------------------- divers
\newlength{\metalab}
\newlength{\tw}\setlength{\tw}{\textwidth}
\newcommand{\settw}[1]{\setlength{\tw}{\dimexpr\textwidth-#1pt\relax}}
% tolérance relâchée : sans elle, TeX préfère déborder de quelques points
% plutôt que de desserrer une ligne contenant un long jeton en chasse fixe.
\emergencystretch=1.6em

% Document autonome : il reprend le préambule de l'équipe (fontes, palette, encadrés) et ajoute
% ce que la chaîne Markdown ne sait pas produire : formules, schémas TikZ et graphiques pgfplots.
% Compilation depuis ce dossier : latexmk -xelatex 03_apprentissage_contrastif.tex

\usepackage{amsmath}
\usepackage{amsthm}
\usepackage{unicode-math}
\setmathfont{LibertinusMath-Regular.otf}
\usepackage{tikz}
\usetikzlibrary{arrows.meta,positioning,calc,fit,backgrounds,babel}
\usepackage{pgfplots}
\pgfplotsset{compat=1.18}

\hypersetup{pdftitle={Entraîner un encodeur de recherche},
  pdfsubject={Fiche pédagogique : recherche dense, évaluation et apprentissage contrastif},
  pdfauthor={Equipe Data Tigers - NLP 2 - M2 Data et IA, FGES}}
\renewcommand{\entetecourant}{Entraîner un encodeur de recherche \textbullet{} Équipe Data Tigers}

\newcolumntype{R}[1]{>{\raggedleft\arraybackslash\hspace{0pt}}p{#1}}
\newcolumntype{Y}{>{\raggedright\arraybackslash\hspace{0pt}}X}
\newcommand{\mesure}{\textcolor{vert}{\faRuler}\,}
\newcommand{\source}{\textcolor{accent}{\faBook}\,}
\newcommand{\vect}[1]{\mathbf{#1}}
\newcommand{\R}{\mathbb{R}}

\tikzset{
  boite/.style={draw=accent,fill=fondclair,rounded corners=2pt,align=center,
    font=\sffamily\small,inner sep=5pt,minimum height=9mm},
  boite2/.style={boite,draw=accent2,fill=white},
  fleche/.style={-{Stealth[length=2.4mm]},draw=encre,line width=0.7pt},
  etiquette/.style={font=\sffamily\scriptsize,text=gris,align=center},
}

\newtheoremstyle{fiche}{6pt}{6pt}{\itshape}{}{\sffamily\bfseries\color{accent}}{.}{0.5em}{}
\theoremstyle{fiche}
\newtheorem{proposition}{Proposition}
\renewcommand{\qedsymbol}{\textcolor{gris}{\faSquare[regular]}}

\addto\captionsfrench{\renewcommand{\contentsname}{Sommaire}}
\begin{document}
\thispagestyle{empty}
\begingroup\sffamily
{\color{accent}\rule{\textwidth}{2.6pt}}\par\vspace{13pt}
{\fontsize{27}{31}\selectfont\bfseries\color{encre}\raggedright Entraîner un encodeur\\de recherche\par}\vspace{7pt}
{\fontsize{13.5}{17}\selectfont\color{gris} Équipe Data Tigers \textbullet{} NLP~2 \textbullet{} M2 Data \& IA \textbullet{} FGES}\par\vspace{5pt}
{\fontsize{11}{14}\selectfont\color{accent2}\bfseries Fiche pédagogique : recherche dense, évaluation et apprentissage contrastif}\par\vspace{12pt}
{\color{filet}\rule{\textwidth}{0.6pt}}\par\vspace{14pt}
\endgroup

\begingroup\setlength{\parskip}{0pt}
\settowidth{\metalab}{\tabfont\bfseries Vérifications}
\begin{tabular}{@{}>{\tabfont\bfseries\color{accent}}l@{\hspace{16pt}}p{\dimexpr\textwidth-\metalab-18pt\relax}@{}}
Objet & Expliquer ce que fait l'encodeur pendant une recherche, ce qui distingue l'évaluation de l'entraînement, et comment l'apprentissage contrastif ajuste ses paramètres \\[3.5pt]
Public & Toute l'équipe ; prolonge le support \emph{Comprendre l'encodeur multilingual-e5-small} (sections 3 à 5), à lire d'abord \\[3.5pt]
Usage & Préparer le notebook 03 (\#46), qui évalue, et l'ajustement fin (\#49), qui entraînerait \\[3.5pt]
Version & 6 octobre 2026 \\[3.5pt]
Vérifications & Chaque chiffre est soit cité d'une source primaire (\source), soit mesuré sur notre corpus (\mesure) \\[3.5pt]
\end{tabular}\par\vspace{16pt}\endgroup

\begin{callout}
Les exemples numériques ont été calculés le 6 octobre 2026 sur les 4\,240 passages du découpage M2, avec le modèle
\code{intfloat/multilingual-e5-small} (révision \code{614241f6}) en précision fp32. Les trois questions employées
sont celles du notebook 02 : aucune ne vient du jeu de test. Elles illustrent un mécanisme ; elles ne mesurent
pas la qualité du système.
\end{callout}

\vspace{10pt}{\color{filet}\rule{\textwidth}{0.4pt}}\par\vspace{16pt}
\begingroup\setlength{\parskip}{0pt}\tableofcontents\endgroup
\clearpage

% =====================================================================================
\section{L'essentiel en cinq points}

\begin{enumerate}
\item \textbf{Le modèle est une fonction} qui transforme \emph{un} texte en \emph{un} vecteur de 384 nombres. Il
  encode chaque texte séparément, question ou passage ; il ne voit jamais un couple d'un seul bloc.
\item \textbf{La matrice de similarités n'est pas le modèle} : c'est un calcul fait à partir de ses vecteurs. Chercher,
  c'est calculer une ligne de cette matrice pour la question posée, puis trier.
\item \textbf{Évaluer n'est pas entraîner.} Les notebooks 01 à 03 et le site utilisent le modèle sans le modifier ; le
  notebook 03 mesure son classement sur nos questions.
\item \textbf{Entraîner, c'est modifier ses 117,7 millions de paramètres} pour que, dans tout lot d'exemples, chaque
  question soit plus proche de son passage que des autres passages du lot. La perte InfoNCE mesure l'écart à cet
  objectif ; la rétropropagation dit comment chaque paramètre doit bouger.
\item \textbf{Le modèle a déjà été entraîné} par ses auteurs. L'ajustement fin (\#49) reprendrait le même mécanisme
  sur des couples de notre domaine, distincts du jeu de test, puis exigerait de recalculer l'index et de remesurer.
\end{enumerate}

% =====================================================================================
\section{Le modèle et ses paramètres}

\subsection{Une fonction des textes vers la sphère}

Notons $\mathcal{T}$ l'ensemble des textes et $d = 384$. L'encodeur est une fonction paramétrée
\[
  f_\theta : \mathcal{T} \longrightarrow \mathbb{S}^{d-1} = \{\vect{x} \in \R^d : \lVert\vect{x}\rVert_2 = 1\},
\]
où $\theta$ désigne l'ensemble de ses paramètres. Pour un texte $x$, le réseau produit un vecteur
$\vect{h}_t \in \R^d$ par token, puis
\[
  \vect{h}(x) = \frac{\sum_t m_t\,\vect{h}_t}{\sum_t m_t},
  \qquad
  f_\theta(x) = \frac{\vect{h}(x)}{\lVert\vect{h}(x)\rVert_2},
\]
avec $m_t = 1$ pour un token réel et $0$ pour un remplissage (support sur l'encodeur, section 3.4). Une question
$q$ et un passage $p$ passent dans le \textbf{même} réseau, distingués par un préfixe \source :
\[
  \vect{u} = f_\theta(\text{\code{query: }} q), \qquad \vect{e} = f_\theta(\text{\code{passage: }} p).
\]

\subsection{Ce que contient \texorpdfstring{$\theta$}{theta}}

\begin{tabletype}
\begin{tabularx}{\textwidth}{@{}YR{2.7cm}R{1.6cm}@{}}
\toprule
\thead{Composant} & \thead{Paramètres} & \thead{Part} \\
\midrule
Table des tokens : 250\,037 lignes de 384 nombres (le tokeniseur en emploie 250\,002) & 96\,014\,208 & 81,6\,\% \\
12 couches Transformer : 12 têtes d'attention, dimension 384, couche intermédiaire de 1\,536 & 21\,293\,568 & 18,1\,\% \\
Positions (512 $\times$ 384), types de segments et normalisation de la table & 198\,144 & 0,2\,\% \\
Couche \emph{pooler}, que la moyenne des tokens n'utilise pas & 147\,840 & 0,1\,\% \\
\midrule
\textbf{Total} & \textbf{117\,653\,760} & 100\,\% \\
\bottomrule
\end{tabularx}
\end{tabletype}

\mesure Décompte fait sur le modèle chargé. Une couche compte $1\,774\,464$ paramètres : $4 \times (384^2 + 384)$
pour l'attention, $2\,(384 \times 1\,536) + 1\,536 + 384$ pour la couche intermédiaire, et $2 \times 768$ pour
ses deux normalisations. Les quatre cinquièmes du modèle servent donc à donner un vecteur de départ à chaque
token du vocabulaire multilingue : c'est ce que vise la réduction du vocabulaire (\#50).

% =====================================================================================
\section{Chercher : une ligne de matrice, puis un tri}

\subsection{Le calcul fait à chaque question}

Le notebook 02 a calculé une fois pour toutes la matrice des passages, ou \textbf{index},
\[
  \vect{E} = \begin{pmatrix} \vect{e}_1^{\top} \\ \vdots \\ \vect{e}_N^{\top} \end{pmatrix} \in \R^{N \times d},
  \qquad N = 4\,240 ,
\]
soit $4\,240 \times 384 \times 4$ octets $= 6{,}21$~Mio en fp32 \mesure. Pour une question $q$ de vecteur $\vect{u}$,
le site calcule
\[
  \vect{v} = \vect{E}\,\vect{u} \in \R^{N}, \qquad v_j = \langle \vect{e}_j, \vect{u} \rangle = \cos(\vect{e}_j, \vect{u}),
\]
puis classe les passages par score décroissant, les ex aequo dans l'ordre de l'index, et affiche les $k$ premiers.
Le calcul de $\vect{v}$ coûte $N d = 1\,628\,160$ multiplications et autant d'additions : 1,7~ms en médiane dans
le navigateur, pour environ 20~ms d'encodage de la question \mesure.

\begin{center}
\begin{tikzpicture}[node distance=5mm and 7mm]
  \node[boite2,text width=2.5cm] (q) {Question\\du visiteur};
  \node[boite,right=of q,text width=2.4cm] (f) {Encodeur $f_\theta$\\\code{query:}};
  \node[boite,right=of f,text width=2.6cm] (v) {$\vect{v} = \vect{E}\vect{u}$\\4\,240 scores};
  \node[boite2,right=of v,text width=2.2cm] (k) {Tri,\\$k$ premiers};
  \node[boite,above=7mm of v,text width=2.6cm] (E) {Index $\vect{E}$\\calculé hors ligne};
  \draw[fleche] (q) -- (f); \draw[fleche] (f) -- node[etiquette,above]{$\vect{u}$} (v);
  \draw[fleche] (v) -- (k); \draw[fleche,dashed] (E) -- (v);
\end{tikzpicture}
\end{center}

\subsection{La matrice de similarités}

Pour évaluer, on pose d'un coup $n$ questions de vecteurs $\vect{u}_1, \dots, \vect{u}_n$, réunis dans
$\vect{U} \in \R^{n \times d}$, et l'on forme
\[
  \vect{S} = \vect{U}\,\vect{E}^{\top} \in \R^{n \times N}, \qquad S_{ij} = \langle \vect{u}_i, \vect{e}_j \rangle .
\]
Chaque ligne est la recherche d'une question. Le notebook 03 construit cette matrice pour les questions du jeu de
test et y lit le rang du premier passage pertinent. Le site, lui, ne la stocke jamais : il calcule à la volée la
seule ligne de la question posée.

\subsection{Pourquoi ne pas tenir une table de scores}

On pourrait imaginer d'apprendre directement les scores, une case par couple (question, passage). Une telle table
n'est définie que pour les questions qu'elle contient : pour une question nouvelle, il n'existe aucune ligne. Or
l'ensemble des questions possibles est illimité, et presque toutes celles du site seront nouvelles. L'encodeur, au
contraire, est défini sur \emph{tout} texte, avec un nombre de paramètres fixé qui ne dépend pas du nombre de
questions : deux formulations voisines reçoivent des vecteurs voisins, ce qui lui permet de \textbf{généraliser}. Une
table apprise serait la forme extrême du surapprentissage : une mémoire, sans généralisation.

\subsection{Bi-encodeur et modèle croisé}

Un modèle \textbf{croisé} lit la question et le passage ensemble, $g_\theta(q, p) \in \R$, et classe souvent mieux.
Mais il ne permet aucun calcul préalable : chaque question exigerait $N = 4\,240$ passages complets dans le
réseau. Au coût d'encodage d'une seule question, environ 20~ms, cela fait déjà plus de 80~secondes par question,
et davantage en réalité, car un couple est bien plus long qu'une question. Le \textbf{bi-encodeur} déplace tout le
travail lourd hors ligne : c'est ce qui rend le volet A possible dans un navigateur (support sur l'encodeur,
section 5.1).

% =====================================================================================
\section{Évaluer n'est pas entraîner}

\begin{tabletype}
\begin{tabularx}{\textwidth}{@{}L{3.0cm}YY@{}}
\toprule
& \thead{Inférence et évaluation} & \thead{Entraînement} \\
\midrule
Paramètres $\theta$ & fixes & modifiés à chaque pas \\
Données & n'importe quel texte ; pour évaluer, des questions annotées & des couples (question, passage qui répond) \\
Résultat & des vecteurs, des classements, des métriques & de nouveaux paramètres $\theta'$ \\
Dans notre projet & notebooks 01, 02 et 03, site du volet A & ajustement fin (\#49), non commencé \\
\bottomrule
\end{tabularx}
\end{tabletype}

\textbf{Évaluer} consiste à lire la matrice $\vect{S}$. Pour la question $i$, soit $r_i$ le rang de son premier
passage pertinent ; avec $n$ questions,
\[
  \text{Recall@}k = \frac{1}{n}\sum_{i=1}^{n} \mathbf{1}[r_i \le k],
  \qquad
  \text{MRR@10} = \frac{1}{n}\sum_{i=1}^{n} \frac{\mathbf{1}[r_i \le 10]}{r_i} .
\]
Ces définitions sont celles de \code{scripts/evaluation/metrics.py}, validées sur PIAF (\#45). Rien n'y est appris.

\textbf{BM25 n'apprend rien non plus.} Sa méthode \code{fit()} compte, pour chaque mot, le nombre de passages qui le
contiennent, et mesure la longueur de chaque passage : ce sont des statistiques, sans gradient. Ses deux réglages,
$k_1$ et $b$, sont fixés à la main. Le nom \code{fit} vient seulement de la convention de scikit-learn.

% =====================================================================================
\section{Entraîner : la matrice d'un lot}

\subsection{Les données et la construction du lot}

On dispose d'un ensemble de couples $\mathcal{D} = \{(q_i, p_i^{+})\}$, où $p_i^{+}$ répond à $q_i$. Un pas
d'entraînement tire un \textbf{lot} de $n$ couples, les encode, et forme la matrice carrée
\[
  \vect{S}^{B} \in \R^{n \times n}, \qquad S_{ij} = \langle \vect{u}_i, \vect{e}_j \rangle,
  \quad \vect{u}_i = f_\theta(\text{\code{query: }} q_i), \quad \vect{e}_j = f_\theta(\text{\code{passage: }} p_j^{+}) .
\]
Les couples étant rangés dans le même ordre en lignes et en colonnes, \textbf{les bons couples sont sur la
diagonale par construction}. Pour la question $i$, les autres colonnes contiennent les passages des autres couples :
ce sont les \textbf{négatifs du lot}, obtenus sans aucune annotation supplémentaire. C'est la stratégie de DPR et
d'E5 \source.

\subsection{Un lot réel sur notre corpus}\label{sec:lot}

\begin{tabletype}
\begin{tabularx}{\textwidth}{@{}YR{1.75cm}R{1.75cm}R{1.75cm}R{1.75cm}@{}}
\toprule
& \thead{$p_1$ mentions interdites} & \thead{$p_2$ arrêt maladie} & \thead{$p_3$ essai, CDI} & \thead{$p_4$ essai, CDD} \\
\midrule
$q_1$ : Quelles mentions sont interdites sur le bulletin de paie ? & \textbf{0,902} & 0,800 & 0,797 & 0,802 \\
$q_2$ : Mon employeur peut-il me licencier pendant un arrêt maladie ? & 0,811 & \textbf{0,885} & 0,805 & 0,808 \\
$q_3$ : Combien de temps peut durer la période d'essai d'un CDI ? & 0,787 & 0,828 & \textbf{0,884} & 0,893 \\
\bottomrule
\end{tabularx}
\end{tabletype}

\mesure Cosinus du modèle actuel. Les colonnes $p_1$ à $p_3$ forment le lot $3 \times 3$ ; $p_4$, la section de la
fiche CDD sur la durée de la période d'essai, est ajouté comme \textbf{négatif difficile} pour $q_3$. C'est l'échec
observé au notebook 02 : $p_4$ (0,893) devance $p_3$ (0,884).

\subsection{Pourquoi les négatifs sont indispensables}

Si l'on cherchait seulement à rapprocher chaque question de son passage, la fonction constante serait une solution
optimale : tous les vecteurs égaux donnent $S_{ii} = 1$, le cosinus maximal, et le modèle ne distingue plus rien.
C'est l'\textbf{effondrement}. Les négatifs rendent l'objectif \emph{relatif} : il ne suffit plus d'être proche de son
passage, il faut l'être \emph{plus} que des autres. La proposition~\ref{prop:effondrement} montre que la perte
InfoNCE pénalise l'effondrement.

% =====================================================================================
\section{La perte InfoNCE}

\subsection{Définition}

Avec une température $\tau > 0$, on pose pour chaque ligne $i$ du lot
\[
  P_{ij} = \frac{\exp(S_{ij}/\tau)}{\sum_{k=1}^{n} \exp(S_{ik}/\tau)},
  \qquad
  \ell_i = -\log P_{ii},
  \qquad
  \mathcal{L} = \frac{1}{n}\sum_{i=1}^{n} \ell_i .
\]
$P_{i\cdot}$ est une loi de probabilité sur les $n$ passages du lot (fonction \emph{softmax}) ; $\ell_i$ est
l'entropie croisée de la classification « parmi ces $n$ passages, lequel est le bon ? ». Des négatifs difficiles
s'ajoutent simplement au dénominateur ; E5 en emploie 7 par exemple lors de son ajustement fin \source. En
réécrivant $\ell_i$,
\[
  \ell_i = \log \sum_{k=1}^{n} \exp\!\Big(\frac{S_{ik} - S_{ii}}{\tau}\Big),
\]
on voit que la perte ne dépend que des \textbf{écarts} entre le score du bon passage et ceux des autres.

\subsection{D'où vient le nom}

L'estimation contrastive par bruit (\emph{Noise-Contrastive Estimation}) apprend un modèle en le faisant distinguer
de vraies données d'échantillons de bruit \source. InfoNCE en est la variante de van den Oord et al. : sous leurs
hypothèses, des négatifs tirés indépendamment, minimiser $\mathcal{L}$ maximise une borne inférieure de
l'information mutuelle entre questions et passages, $I(q; p) \ge \log n - \mathcal{L}$ \source. D'où l'intérêt de
lots nombreux : E5 a été pré-entraîné avec des lots de 32\,768 couples \source.

\subsection{Bornes}

\begin{proposition}\label{prop:bornes}
Pour des vecteurs normalisés, donc des scores dans $[-1, 1]$, chaque terme vérifie
$\ell_i \ge \log\big(1 + (n-1)\,e^{-2/\tau}\big) > 0$.
\end{proposition}

\begin{proof}
$P_{ii}$ est maximal quand $S_{ii} = 1$ et $S_{ik} = -1$ pour $k \neq i$ ; alors
$P_{ii} = e^{1/\tau} / \big(e^{1/\tau} + (n-1)e^{-1/\tau}\big) = 1 / \big(1 + (n-1)e^{-2/\tau}\big)$, et
$\ell_i = -\log P_{ii}$.
\end{proof}

Avec $\tau = 0{,}01$ et $n = 32$, cette borne vaut $\log(1 + 31\,e^{-200})$, soit pratiquement zéro : la perte peut
devenir aussi petite que l'on veut, sans jamais s'annuler.

\begin{proposition}\label{prop:effondrement}
Si $f_\theta$ est constante, alors $\mathcal{L} = \log n$, quelle que soit $\tau$.
\end{proposition}

\begin{proof}
Tous les vecteurs sont égaux, donc $S_{ik} = 1$ pour tous $i, k$, puis $P_{ii} = 1/n$ et $\ell_i = \log n$.
\end{proof}

Pour $n = 32$, l'effondrement coûte $\log 32 \approx 3{,}466$, très loin du minimum : il n'est donc pas une
solution de la minimisation.

\subsection{Le rôle de la température}

Comme $\ell_i = \log \sum_k \exp\big((S_{ik} - S_{ii})/\tau\big)$, deux limites éclairent son rôle :
\begin{itemize}
\item quand $\tau \to 0$, la somme est dominée par son plus grand terme, et
  $\ell_i \approx \max\big(0,\ \max_{k \neq i} S_{ik} - S_{ii}\big)/\tau$ : seule compte la question de savoir si le
  plus dangereux des négatifs passe devant le bon passage, et de combien ;
\item quand $\tau \to \infty$, tous les termes tendent vers 1, $P_{i\cdot}$ devient uniforme et $\ell_i \to \log n$ :
  le signal d'apprentissage disparaît.
\end{itemize}
E5 emploie $\tau = 0{,}01$ \source. C'est cohérent avec ses scores resserrés : deux passages pris au hasard sur notre
corpus ont un cosinus médian de 0,834 \mesure (support sur l'encodeur, section 4.4). Sans une petite température, des
écarts de quelques centièmes resteraient sans effet.

\begin{tabletype}
\begin{tabularx}{\textwidth}{@{}YR{2.7cm}R{2.7cm}@{}}
\toprule
\thead{Lot réel de la section \ref{sec:lot}} & \thead{$\tau = 0{,}01$ (E5)} & \thead{$\tau = 0{,}05$} \\
\midrule
Lot $3 \times 3$ : probabilité du bon passage pour $q_1$, $q_2$, $q_3$ & $>0{,}999$ ; 0,999 ; 0,996 & 0,798 ; 0,698 ; 0,680 \\
Lot $3 \times 3$ : perte $\mathcal{L}$ & 0,002 & 0,324 \\
Avec le négatif difficile $p_4$ : probabilité du bon passage pour $q_3$ & 0,290 & 0,375 \\
Avec le négatif difficile $p_4$ : perte $\ell_3$ & 1,237 & 0,981 \\
\bottomrule
\end{tabularx}
\end{tabletype}

\mesure Le lot $3 \times 3$ est déjà résolu par le modèle actuel : avec $\tau = 0{,}01$, sa perte est presque nulle,
et un tel lot n'apprendrait presque rien. Le négatif difficile change tout : $q_3$ ne place plus que 29\,\% de sa
probabilité sur le bon passage, et sa perte passe de 0,004 à 1,237.

\begin{center}
\begin{tikzpicture}
\begin{axis}[width=0.86\textwidth,height=5.6cm,xmode=log,xmin=0.002,xmax=1,ymin=0,ymax=4.7,
  xlabel={température $\tau$},ylabel={perte $\ell_3$},
  tick label style={font=\sffamily\scriptsize},label style={font=\sffamily\small},
  legend style={font=\sffamily\scriptsize,draw=none,fill=none,at={(0.98,0.97)},anchor=north east},
  grid=major,grid style={filet}]
  \addplot[accent2,line width=1pt,domain=-6.2:0,samples=160]
    ({exp(x)},{ln(exp(-0.096719/exp(x))+exp(-0.055906/exp(x))+1+exp(0.008931/exp(x)))});
  \addlegendentry{avec le négatif difficile $p_4$ (4 passages)}
  \addplot[accent,line width=1pt,domain=-6.2:0,samples=160]
    ({exp(x)},{ln(exp(-0.096719/exp(x))+exp(-0.055906/exp(x))+1)});
  \addlegendentry{lot $3 \times 3$ seul}
  \draw[gris,dashed] (axis cs:0.01,0) -- (axis cs:0.01,4.7) node[pos=0.93,right,font=\sffamily\scriptsize]{E5};
  \addplot[gris,dotted,domain=0.002:1] {ln(4)};
  \addplot[gris,dotted,domain=0.002:1] {ln(3)};
\end{axis}
\end{tikzpicture}
\end{center}

\mesure Perte de $q_3$ en fonction de $\tau$, calculée à partir des scores exacts du lot de la section \ref{sec:lot}. Les
pointillés marquent $\log 4$ et $\log 3$, vers lesquelles chaque courbe tend quand $\tau$ grandit. Quand $\tau$
diminue, la perte s'annule si le bon passage est en tête (courbe bleue), et croît comme $0{,}009/\tau$ si un négatif
le devance de 0,009 (courbe brune).

% =====================================================================================
\section{Le gradient : qui pousse qui}

\subsection{Par rapport aux scores}

\begin{proposition}\label{prop:gradient}
$\displaystyle \frac{\partial \ell_i}{\partial S_{ij}} = \frac{1}{\tau}\big(P_{ij} - \mathbf{1}[i = j]\big)$.
\end{proposition}

\begin{proof}
$\ell_i = -S_{ii}/\tau + \log \sum_k \exp(S_{ik}/\tau)$. La dérivée du premier terme vaut $-\mathbf{1}[i=j]/\tau$ ;
celle du second vaut $\exp(S_{ij}/\tau)/\big(\tau\sum_k \exp(S_{ik}/\tau)\big) = P_{ij}/\tau$.
\end{proof}

La descente de gradient déplace chaque score dans le sens opposé à sa dérivée. Le bon score $S_{ii}$ est donc
\textbf{tiré vers le haut} avec un poids $(1 - P_{ii})/\tau$, et chaque négatif \textbf{repoussé} avec un poids
$P_{ij}/\tau$. Comme $\sum_j P_{ij} = 1$, la traction exercée sur le bon passage est exactement répartie sur les
négatifs, \textbf{au prorata de leur probabilité}. Un négatif facile ($P_{ij} \approx 0$) ne reçoit presque rien ; un
négatif difficile reçoit presque tout.

\mesure Sur le lot réel avec $\tau = 0{,}01$, la ligne de $q_3$ vaut
$P_{3\cdot} = (0{,}000 ;\ 0{,}001 ;\ 0{,}290 ;\ 0{,}709)$. Le bon passage $p_3$ est tiré avec le poids
$1 - 0{,}290 = 0{,}710$ ; le passage CDD en reçoit 0,709, soit la quasi-totalité. Les passages sur le bulletin
de paie et la maladie, faciles, ne reçoivent pratiquement rien : un lot de négatifs aléatoires apprend peu.

\subsection{Jusqu'aux vecteurs}

Comme $S_{ij} = \langle \vect{u}_i, \vect{e}_j \rangle$, on a $\partial S_{ij}/\partial \vect{u}_i = \vect{e}_j$ et
$\partial S_{ij}/\partial \vect{e}_j = \vect{u}_i$, donc
\[
  \frac{\partial \ell_i}{\partial \vect{u}_i} = \frac{1}{\tau}\Big(\sum_{j} P_{ij}\,\vect{e}_j - \vect{e}_i\Big).
\]
Un pas de gradient déplace $\vect{u}_i$ vers $\vect{e}_i$ et l'éloigne de la moyenne des passages pondérée par
$P_{i\cdot}$, c'est-à-dire surtout des négatifs dangereux. Symétriquement, chaque $\vect{e}_j$ bouge.

Ces vecteurs restent sur la sphère, car ils sont normalisés. Avec $f = \vect{h}/\lVert\vect{h}\rVert$, la dérivée de la
normalisation vaut
\[
  \frac{\partial f}{\partial \vect{h}} = \frac{1}{\lVert\vect{h}\rVert}\big(\vect{I} - f f^{\top}\big),
\]
qui projette le gradient sur le plan tangent à la sphère : l'apprentissage fait \emph{tourner} les vecteurs, il ne
les allonge pas.

\subsection{Jusqu'aux paramètres : la rétropropagation}

Les vecteurs dépendent de $\theta$ à travers la normalisation, la moyenne des tokens
($\partial \vect{h}/\partial \vect{h}_t = m_t / \sum_s m_s$), les 12 couches et la table des tokens. La
\textbf{rétropropagation} applique la règle de dérivation en chaîne de la perte vers l'entrée, couche par couche, et
fournit en une seule passe arrière la dérivée de $\mathcal{L}$ par rapport à chacun des 117\,653\,760 paramètres.
Comme questions et passages passent dans le même réseau, les deux côtés contribuent au gradient des mêmes
paramètres : un pas déplace à la fois les vecteurs des questions et ceux des passages.

% =====================================================================================
\section{Optimiser : la descente de gradient}

\subsection{Le pas}

La descente de gradient stochastique remplace $\theta$ par
\[
  \theta_{t+1} = \theta_t - \eta\, \vect{g}_t, \qquad \vect{g}_t = \nabla_\theta \mathcal{L}(\theta_t ; B_t),
\]
où $B_t$ est le lot du pas $t$ et $\eta$ le pas d'apprentissage. Le gradient d'un lot est une estimation bruitée du
gradient sur toutes les données ; c'est le sens de « stochastique ». En pratique, on emploie AdamW, qui adapte le
pas de chaque paramètre à l'aide de moyennes mobiles des gradients et de leurs carrés, et applique une décroissance
des poids découplée \source.

\subsection{Les ordres de grandeur pour nous}

\mesure Sur la machine de référence, un pas sur un lot de 32 couples prend 89~ms (journal d'usage de l'IA,
session 8). Avec $17\,106$ couples, une époque compte $\lceil 17\,106 / 32 \rceil = 535$ pas, et trois époques
$1\,605$ pas, soit environ 143~secondes : 2,4~minutes. Le coût de calcul n'est pas un obstacle. Avec des lots de
32, chaque question n'a que 31 négatifs, contre 32\,767 pour E5 ; les auteurs notent que des négatifs difficiles
permettent de compenser des lots plus petits \source.

\subsection{Deux risques}

\begin{itemize}
\item Le \textbf{surapprentissage} : à force de pas, le modèle apprend par cœur les couples d'entraînement et
  généralise moins bien. Parades : un petit pas $\eta$, peu d'époques, et l'arrêt quand la perte d'un ensemble de
  validation, distinct de l'entraînement et du test, cesse de baisser.
\item L'\textbf{oubli catastrophique} : en s'adaptant à notre domaine, le modèle peut perdre une partie de ce qu'il
  savait du français général. Seule une mesure avant et après, sur le même jeu figé, permet de le voir.
\end{itemize}

% =====================================================================================
\needspace{12\baselineskip}
\section{Ce que cela implique pour notre projet}

\subsection{Où nous en sommes}

Le modèle a été entraîné par ses auteurs : pré-entraînement contrastif sur environ un milliard de couples de textes,
puis ajustement fin sur environ 1,6 million de couples annotés, avec négatifs difficiles et distillation depuis un
modèle croisé \source. Nous l'utilisons tel quel. Le notebook 03 fournira la mesure \textbf{avant} tout ajustement.

\subsection{Une recette pour l'ajustement fin (\#49)}

Proposition, à décider dans le ticket :
\begin{enumerate}
\item \textbf{Données.} Des couples (question, passage qui répond) construits sur les passages M2, par exemple avec
  des questions générées par IA, que le README du jeu de questions réserve à l'entraînement. \textbf{Jamais une
  question du jeu de test}, sans quoi la mesure après ajustement évaluerait de la mémoire.
\item \textbf{Validation.} Un ensemble de validation tiré des couples d'entraînement, séparé par fiche, pour
  arrêter l'entraînement à temps.
\item \textbf{Lots sans faux négatifs.} Des lots composés de fiches différentes : deux couples de la même fiche
  feraient d'un passage pertinent un négatif. Le découpage M2 compte 57 passages à texte identique \mesure (support
  sur l'encodeur, section 8.4).
\item \textbf{Négatifs difficiles.} Pour chaque question, des passages que le modèle actuel classe haut sans qu'ils
  répondent, vérifiés : le passage CDD pour une question sur le CDI en est l'exemple type.
\item \textbf{Après l'entraînement.} $f_\theta$ a changé : recalculer l'index des 4\,240 passages (notebook 02),
  réexporter le modèle en ONNX q8 et contrôler sa fidélité dans le navigateur, puis remesurer avec le notebook 03,
  selon le même protocole et sur le même jeu figé.
\end{enumerate}

\begin{center}
\begin{tikzpicture}[node distance=5mm and 8mm]
  \node[boite2,text width=3.0cm] (train) {Couples d'entraînement\\(questions générées)};
  \node[boite,right=of train,text width=2.6cm] (ft) {Ajustement fin\\(\#49) : $\theta \to \theta'$};
  \node[boite,right=of ft,text width=2.6cm] (idx) {Nouvel index\\(notebook 02)};
  \node[boite,right=of idx,text width=2.4cm] (eval) {Évaluation\\(notebook 03)};
  \node[boite2,above=8mm of eval,text width=2.4cm] (test) {Jeu de test figé\\(\#44)};
  \draw[fleche] (train) -- (ft); \draw[fleche] (ft) -- (idx); \draw[fleche] (idx) -- (eval);
  \draw[fleche] (test) -- (eval);
  \draw[rouge,line width=0.8pt,dashed,-{Stealth[length=2.4mm]}] (test.west) -| (train.north)
    node[pos=0.3,above,etiquette,text=rouge]{jamais vers l'entraînement}
    node[pos=0.5,fill=white,inner sep=1pt,text=rouge]{\faTimes};
\end{tikzpicture}
\end{center}

% =====================================================================================
\section{Glossaire}

\begin{description}[leftmargin=0pt,labelsep=0.5em,itemsep=2pt]
\item[Encodeur] fonction $f_\theta$ qui transforme un texte en vecteur normalisé.
\item[Paramètres] les nombres $\theta$ du réseau, 117\,653\,760 ici, que l'entraînement modifie.
\item[Inférence] utilisation du modèle à paramètres fixes : encoder, chercher, évaluer.
\item[Lot] ensemble de $n$ couples traités ensemble à un pas d'entraînement.
\item[Négatif du lot] passage d'un autre couple du lot, employé comme mauvaise réponse.
\item[Négatif difficile] passage proche de la question sans y répondre, choisi exprès.
\item[Faux négatif] passage traité comme négatif alors qu'il répond à la question.
\item[Softmax] transformation de scores en probabilités : $P_{ij} \propto \exp(S_{ij}/\tau)$.
\item[Température] constante $\tau$ qui règle la sévérité de la comparaison des scores.
\item[InfoNCE] perte $-\log P_{ii}$ moyennée sur le lot : entropie croisée de la classification du bon passage.
\item[Effondrement] solution triviale où tous les textes reçoivent le même vecteur.
\item[Gradient] vecteur des dérivées de la perte par rapport aux paramètres.
\item[Rétropropagation] calcul de ce gradient par dérivation en chaîne, de la perte vers l'entrée.
\item[Époque] passage complet sur les données d'entraînement.
\item[Surapprentissage] apprentissage par cœur des exemples, au détriment de la généralisation.
\item[Oubli catastrophique] perte de compétences générales lors d'une adaptation à un domaine.
\item[Bi-encodeur] question et passage encodés séparément puis comparés ; passages encodés à l'avance.
\item[Modèle croisé] question et passage lus ensemble ; plus précis, sans calcul préalable possible.
\end{description}

% =====================================================================================
\section{Sources}

\begin{itemize}
\item Gutmann M. et Hyvärinen A. (2010), {\NoAutoSpacing\emph{Noise-contrastive estimation: A new estimation principle for
  unnormalized statistical models}}, AISTATS.
\item van den Oord A., Li Y. et Vinyals O. (2018), \emph{Representation Learning with Contrastive Predictive
  Coding}, {\NoAutoSpacing arXiv:1807.03748}.
\item Karpukhin V. et al. (2020), \emph{Dense Passage Retrieval for Open-Domain Question Answering}, EMNLP :
  négatifs du lot.
\item Loshchilov I. et Hutter F. (2019), \emph{Decoupled Weight Decay Regularization}, ICLR : AdamW.
\item Wang L. et al. (2022), \emph{Text Embeddings by Weakly-Supervised Contrastive Pre-training},
  {\NoAutoSpacing arXiv:2212.03533}, sections 4.1 et 4.2 et section 5 : perte, $\tau = 0{,}01$, préfixes, négatifs du lot, lots de
  32\,768, 7 négatifs difficiles.
\item Wang L. et al. (2024), {\NoAutoSpacing\emph{Multilingual E5 Text Embeddings: A Technical Report}, arXiv:2402.05672},
  tableaux 1 et 2 : données de pré-entraînement et d'ajustement fin.
\item Équipe Data Tigers : support \emph{Comprendre l'encodeur multilingual-e5-small} ; rapport de validation des
  métriques sur PIAF ; journal d'usage de l'IA, session 8.
\end{itemize}

\end{document}
